% !PRINTED 221
\weberpage{221}

The computation of the coefficients $Q$ is easy to carry out if one sets up this equation for $x_1$, $x_2$, $x_3$ and correspondingly $y_1$, $y_2$, $y_3$.

We shall give a few indications of the course of the computation, which requires only the representation of symmetric functions of the roots of a cubic equation by the coefficients, according to the prescriptions of the fourth section (\S.\,42, 45). If one first takes the sum of the three equations (9), one obtains by (6)
\[
3 Q_0 = \frac{2}{3} H(t_1, t_0) Q_2 - a_1 f(t_1, t_0),
\]
so that only $Q_1$ and $Q_2$ still need to be computed; these one finds by multiplying the equation (9) formed for $x = x_1$ by
\[
y_2 - y_3 = a_0 (x_2 - x_3) (t_1 - t_0 x_1)
\]
and by
\[
y_2^2 - y_3^2 = -a_0 (x_2 - x_3) y_1 (t_1 - t_0 x_1)
\]
and adding the three analogous equations formed by cyclic permutation of the indices 1, 2, 3.

One then has only to make use of the two formulae
\[
a_0^2 \Sigma x_1^2 (x_2 - x_3) = -\sqrt{D}
\]
\[
a_0^3 \Sigma x_1^3 (x_2 - x_3) = a_1 \sqrt{D},
\]
and thus finds
\[
Q_0 = \frac{-a_1}{3} f(t_1, t_0) - \frac{2}{3} t_0 H(t_1, t_0)
\]
\[
Q_1 = a_0 t_1^2 + \frac{2}{3} a_1 t_0 t_1 + \frac{1}{3} a_2 t_0^2
\]
\[
Q_2 = -t_0.
\]

If for simplification we set $t_1 = t$, $t_0 = 1$, thus passing to the inhomogeneous expressions, we obtain
\weq{10}{
(3 a_0 x + a_1) f(t) = -\frac{2}{3} H(t) + y f'(t) - 3 y^2,
}
while the equation for $y$ reads
\weq{11}{
y^3 + \frac{1}{3} H(t) y + \frac{1}{27} Q(t) = 0,
}
if we write $H(t)$, $Q(t)$ for $H(t, 1)$, $Q(t, 1)$.

Now, in order to solve the cubic equation, we determine $t$ from the quadratic equation
\[
H(t) = 0.
\]
